\section{Contribución 1: Intent-Conditioned Controllable Dialogue}

Si se deja que el modelo de lenguaje continúe libremente a partir del estado completo, tiene que decidir a la vez «qué debe hacer» y «cómo debe decirlo», y además puede filtrar planes privados o inventar acciones ilegales. Cicero primero deja que el módulo estratégico decida la intención de corto plazo y después permite que el modelo de lenguaje se concentre en expresar esa intención como un mensaje natural y oportuno.

\subsection{Por qué se necesita una intent explícita}

El primer punto de la slide con el panorama de las contribuciones es el controllable dialogue: el modelo no solo hace condition on el game state y el dialogue history, sino también condition on las intended actions. Esta interfaz separa la «corrección estratégica» y la «realización lingüística» en dos capas que pueden evaluarse por separado.

\lecturefigure{slide-25-main-contribution-dialogue.jpg}{Main contribution: controlar la generación de diálogo con intended actions}{Diapositiva V025 recuperada de la grabación oficial; Stanford CS25 Lecture 14}

\begin{knowledgebox}{Aquí la intent no es un estado mental abstracto}
La intent de Cicero consiste principalmente en acciones ejecutables de corto plazo, por ejemplo, hacia dónde se prepara una unidad para moverse, apoyar o defender en este turno o en el siguiente. Está más estructurada que una oración en lenguaje natural, lo que facilita verificar que sea coherente con el board state; pero también es más estrecha que un objetivo estratégico de largo plazo, una limitación a la que se volverá más adelante.
\end{knowledgebox}

\subsection{Cómo construir etiquetas de intent a partir de diálogos humanos}

Los registros humanos no tienen anotado manualmente «el plan detrás de esta frase». El sistema primero entrena un intent model que, a partir del state, el history y el diálogo, predice las actions del hablante y del receptor al final del turno; después usa las orders reales para generar automáticamente etiquetas de intención de corto plazo para los mensajes históricos. Para reducir el riesgo de tomar mentiras como objetivo de entrenamiento, el proceso filtra los turns con alta probabilidad de no ser veraces.

\lecturefigure{slide-26-intent-model-training.jpg}{Intent model training: inferir, a partir del diálogo y las acciones finales, el plan correspondiente a cada mensaje}{Diapositiva V026 recuperada de la grabación oficial; Bakhtin et al., Science 2022}

\lecturefigure{slide-27-intent-annotation.jpg}{Intent annotation: agregar automáticamente a los mensajes humanos las intenciones de acción del speaker y del recipient}{Diapositiva V027 recuperada de la grabación oficial; Stanford CS25 Lecture 14}

\teachervoice{El docente aclara en particular que las etiquetas de entrenamiento provienen de turns relativamente confiables, en lugar de suponer que todos los mensajes humanos son veraces. Esta decisión sacrifica parte de la diversidad lingüística, pero vuelve más claro el problema de controllability: primero se aprende a expresar de forma estable el plan real y después se aborda el engaño, que es más complejo.}

\subsection{Cómo se usa la intent durante la inferencia}

Durante la ejecución en línea, el planner primero propone las acciones que Cicero quiere que tomen él mismo y su interlocutor; el intent model las organiza como condiciones; luego el dialogue model, junto con el state y el history, genera varios mensajes candidatos. Así, el modelo de lenguaje no tiene que reinventar el objetivo entre una enorme cantidad de estrategias posibles, sino que resuelve un problema de generación condicional más controlable.

\lecturefigure{slide-28-dialogue-inference.jpg}{Dialogue inference: el resultado de la planeación se convierte primero en intent y luego restringe los mensajes candidatos}{Diapositiva V028 recuperada de la grabación oficial; Stanford CS25 Lecture 14}

Si $m$ es el mensaje por enviar, y $z_s$ y $z_r$ representan, respectivamente, las intenciones de acción de corto plazo del speaker y del recipient, puede escribirse
\begin{equation}
p(m\mid s,h,d,z_s,z_r).
\end{equation}
A diferencia de modelar solo $p(m\mid d)$, esta distribución condicional exige de forma explícita que el mensaje sea coherente a la vez con la situación del tablero, el historial y el plan. No garantiza que el contenido generado sea veraz o eficaz, pero ofrece un punto de apoyo estructurado para las verificaciones de coherencia.

El ejemplo de mapa de la slide 29 muestra la intent y el message lado a lado. Al leer la figura, primero hay que observar las flechas de acción y después cómo el lenguaje convierte las acciones en solicitudes o compromisos: un mismo conjunto de orders admite redacciones distintas, pero si un mensaje pide un apoyo que contradice las flechas, debe clasificarse como plan-inconsistent.

\lecturefigure{slide-29-intent-conditioned-dialogue.jpg}{Intent-conditioned dialogue: de las intenciones de acción en el mapa a solicitudes en lenguaje natural}{Diapositiva V029 recuperada de la grabación oficial; Stanford CS25 Lecture 14}

\subsection{Resultados: por qué mejoran a la vez la controlabilidad y la calidad}

La slide 30 compara tres entradas: el language model puro, el que agrega game-state grounding y el que además agrega intent grounding. Las cuatro columnas son coherencia con el state, coherencia con el plan, proporción de mensajes high-quality según evaluación humana y perplexity. El Cicero completo alcanza 87.30\%, 92.86\%, 37.30\% y 7.70, y supera en las cuatro métricas a las dos configuraciones anteriores.

\lecturefigure{slide-30-intent-quality-results.jpg}{El intent grounding mejora a la vez la state consistency, la plan consistency, la calidad y la perplexity}{Diapositiva V030 recuperada de la grabación oficial; Bakhtin et al., Science 2022}

\begin{knowledgebox}{Primera explicación de la perplexity}
La perplexity es la exponencial de la cross-entropy a nivel de token: $\mathrm{PPL}=\exp(H)$. Intuitivamente, indica cuántas «opciones equivalentes» enfrenta el modelo en cada paso; un valor más bajo suele indicar que la asignación de probabilidad sobre el held-out human dialogue está más concentrada. Pero la perplexity no es un beneficio estratégico: debe interpretarse junto con la coherencia, la calidad según evaluación humana y los resultados reales de las partidas.
\end{knowledgebox}

\begin{importantbox}{Por qué agregar restricciones mejora el lenguaje}
El intent conditioning acota el problema de generación: el modelo no necesita inventar una estrategia mientras organiza el lenguaje, sino solo elegir una expresión adecuada para un objetivo dado. Por eso, la restricción no es solo una barrera de seguridad, sino también una forma de descomposición de la tarea que reduce el nonsense y permite que el modelo de lenguaje aproveche mejor su capacidad limitada.
\end{importantbox}

\subsection{Resumen de la sección}

El intent-conditioned dialogue separa la decisión estratégica de su realización lingüística superficial. La anotación automática proporciona la señal de entrenamiento, la intent estructurada proporciona la interfaz de control y la tabla de resultados muestra que esta descomposición no solo aumenta la plan consistency, sino que también mejora la calidad general del lenguaje.
